% =====================================================================
%  HKA Weekly-Präsentation – AI Lab (WS 26/27)
%  Gruppe GlEr · "Can we detect AI-written text?"
%
%  Stil angelehnt an die Kick-off-Folien (16:9, heller Hintergrund,
%  HKA-Blau für Titel, HKA-Rot für Kicker/Akzente).
%
%  Kompilieren:  latexmk -pdf weekly_template.tex
%                (oder 2x pdflatex)
%
%  Optional: Offizielles HKA-Logo als "hka-logo.pdf" (oder .png) in
%  diesen Ordner legen – sonst wird ein Text-Platzhalter gezeigt.
% =====================================================================
\documentclass[aspectratio=169,11pt]{beamer}

% ---------------------------------------------------------------------
%  Pakete
% ---------------------------------------------------------------------
\usepackage[T1]{fontenc}
\usepackage[utf8]{inputenc}
\usepackage[english]{babel}        % für deutsche Folien: [ngerman]
\usepackage[sfdefault]{plex-sans}  % IBM Plex Sans (wie Originalfolien)
\usepackage{plex-mono}
\usepackage{graphicx}
\usepackage{booktabs}
\usepackage{tabularx}
\usepackage{array}
\usepackage{colortbl}
\usepackage{xcolor}
\usepackage{tikz}
\usepackage{listings}

% ---------------------------------------------------------------------
%  Metadaten – hier jede Woche anpassen
% ---------------------------------------------------------------------
\newcommand{\course}{AI Lab}
\newcommand{\week}{Week 1}
\newcommand{\groupname}{Group GlEr}
\newcommand{\projecttopic}{Can we detect AI-written text?}

\title{\course}
\subtitle{\projecttopic}
\author{Vorname Nachname \and Vorname Nachname}
\institute{\groupname}
\date{Weekly Update \textperiodcentered{} \week{} \textperiodcentered{} \today}

% Fußzeile links (wird auf jeder Inhaltsfolie gezeigt)
\newcommand{\footertext}{\course{} \textperiodcentered{} \groupname{} \textperiodcentered{} \week}

% ---------------------------------------------------------------------
%  Farben (aus den HKA-Kick-off-Folien abgegriffen)
% ---------------------------------------------------------------------
\definecolor{hkaBlue}{HTML}{114F8E}
\definecolor{hkaBlueTitle}{HTML}{0E5195}
\definecolor{hkaRed}{HTML}{C7201A}
\definecolor{hkaText}{HTML}{1E2936}
\definecolor{hkaGray}{HTML}{495467}
\definecolor{hkaLight}{HTML}{8A93A3}
\definecolor{hkaRule}{HTML}{E1E6EB}
\definecolor{hkaStripe}{HTML}{EEF0F1}
\definecolor{hkaBg}{HTML}{FAFAF7}

% ---------------------------------------------------------------------
%  Beamer-Theme
% ---------------------------------------------------------------------
\usetheme{default}
\setbeamersize{text margin left=12mm, text margin right=12mm}

\setbeamercolor{background canvas}{bg=hkaBg}
\setbeamercolor{normal text}{fg=hkaText}
\setbeamercolor{structure}{fg=hkaBlue}
\setbeamercolor{alerted text}{fg=hkaRed}
\setbeamercolor{example text}{fg=hkaBlue}
\setbeamercolor{frametitle}{fg=hkaBlue}
\setbeamercolor{framesubtitle}{fg=hkaGray}
\setbeamercolor{title}{fg=hkaBlueTitle}
\setbeamercolor{subtitle}{fg=black}
\setbeamercolor{block title}{fg=hkaBlue,bg=hkaStripe}
\setbeamercolor{block body}{fg=hkaText,bg=hkaStripe!50}
\setbeamercolor{block title alerted}{fg=white,bg=hkaRed}
\setbeamercolor{block body alerted}{fg=hkaText,bg=hkaRed!6}
\setbeamercolor{block title example}{fg=white,bg=hkaBlue}
\setbeamercolor{block body example}{fg=hkaText,bg=hkaBlue!6}

\setbeamerfont{frametitle}{size=\LARGE,series=\bfseries}
\setbeamerfont{framesubtitle}{size=\normalsize,series=\mdseries}
\setbeamerfont{title}{size=\fontsize{38}{44}\selectfont,series=\bfseries}
\setbeamerfont{subtitle}{size=\Large,series=\mdseries}
\setbeamerfont{block title}{series=\bfseries}

\setbeamertemplate{navigation symbols}{}
\setbeamertemplate{blocks}[rounded][shadow=false]
\setbeamertemplate{itemize items}{\raisebox{0.15ex}{\small\textcolor{hkaRed}{\textbullet}}}
\setbeamertemplate{itemize subitem}{\textcolor{hkaGray}{--}}
\setbeamertemplate{enumerate items}[default]
\setbeamercolor{enumerate item}{fg=hkaRed}
\setlength{\leftmargini}{1.2em}

% --- Kicker: kleine rote Zeile über dem Folientitel ------------------
%  Standard = aktueller \section-Name. Pro Folie überschreibbar mit
%  \kicker{...} direkt VOR \begin{frame}.
\makeatletter
\newcommand{\hka@kicker}{}
\newcount\hka@kickerframe
\newcommand{\kicker}[1]{\gdef\hka@kicker{#1}%
  \global\hka@kickerframe=\numexpr\c@framenumber+1\relax}
\newcommand{\hka@printkicker}{%
  \ifnum\c@framenumber=\hka@kickerframe\else\gdef\hka@kicker{}\fi
  \ifx\hka@kicker\@empty
    \def\hka@k{\insertsectionhead}%
  \else
    \let\hka@k\hka@kicker
  \fi
  \setbox0\hbox{\hka@k}%
  \ifdim\wd0>0pt
    {\usebeamercolor[fg]{alerted text}\footnotesize\bfseries
     \addfontfeature{LetterSpace=6}\MakeUppercase{\hka@k}}\par\vspace{1.2mm}%
  \fi
}
\makeatother
% plex-sans unter pdflatex kennt kein \addfontfeature → Dummy
\providecommand{\addfontfeature}[1]{}

\makeatletter
\setbeamertemplate{frametitle}{%
  \vspace*{4mm}%
  \begin{minipage}[t]{\textwidth}
    \hka@printkicker
    {\usebeamerfont{frametitle}\usebeamercolor[fg]{frametitle}\insertframetitle\par}%
    \ifx\insertframesubtitle\@empty\else
      \vspace{1.5mm}%
      {\usebeamerfont{framesubtitle}\usebeamercolor[fg]{framesubtitle}\insertframesubtitle\par}%
    \fi
  \end{minipage}%
  \vspace{3mm}%
}
\makeatother

% --- Fußzeile: Text links, Seitenzahl rechts ------------------------
\setbeamertemplate{footline}{%
  \leavevmode
  \hbox to \paperwidth{%
    \hspace{12mm}%
    {\scriptsize\color{hkaLight}\footertext}%
    \hfill
    {\scriptsize\color{hkaLight}\insertframenumber}%
    \hspace{12mm}}%
  \vspace{5mm}%
}

% --- HKA-Logo (Datei oder Text-Platzhalter) --------------------------
\newcommand{\hkalogo}[1][9mm]{%
  \IfFileExists{hka-logo.pdf}{\includegraphics[height=#1]{hka-logo.pdf}}{%
  \IfFileExists{hka-logo.png}{\includegraphics[height=#1]{hka-logo.png}}{%
    \begin{tikzpicture}[baseline=(h.base)]
      \node[anchor=east, align=left, inner sep=0, text=hkaRed,
            font=\fontsize{4.5}{5.2}\selectfont] (t)
        {\textbf{Hochschule Karlsruhe}\\University of\\Applied Sciences};
      \node[anchor=west, inner sep=0, text=hkaRed,
            font=\fontsize{20}{20}\selectfont\bfseries] (h)
        at ([xshift=3mm]t.east) {HKA};
    \end{tikzpicture}}}%
}

% --- Titelfolie (zentriert wie im Original) -------------------------
\setbeamertemplate{title page}{%
  \vfill
  \centering
  {\usebeamerfont{title}\usebeamercolor[fg]{title}\inserttitle\par}
  \vspace{4mm}
  {\usebeamerfont{subtitle}\usebeamercolor[fg]{subtitle}\insertsubtitle\par}
  \vspace{5mm}
  \hkalogo\par
  \vspace{5mm}
  {\normalsize\color{black}\insertauthor\par}
  \vspace{1.5mm}
  {\small\color{hkaGray}\insertinstitute\par}
  \vspace{1.5mm}
  {\small\color{black}\insertdate\par}
  \vfill
}

% --- Abschnitts-Trennfolie automatisch bei \section -----------------
\AtBeginSection[]{%
  \begingroup
  \setbeamertemplate{footline}{}
  \begin{frame}[noframenumbering]
    \vfill
    {\usebeamercolor[fg]{alerted text}\footnotesize\bfseries
     \MakeUppercase{\week}\par}
    \vspace{2mm}
    {\fontsize{30}{36}\selectfont\bfseries\color{hkaBlue}\insertsectionhead\par}
    \vspace{4mm}
    {\color{hkaRed}\rule{25mm}{1.2pt}}
    \vfill
  \end{frame}
  \endgroup
}

% --- Schlüssel/Wert-Tabelle wie auf den Kick-off-Folien -------------
%  \begin{kvtable}
%    \kv{Dataset}{...}
%    \kvred{Experiment 1}{...}   % rote Beschriftung
%  \end{kvtable}
\newenvironment{kvtable}{%
  \small
  \renewcommand{\arraystretch}{1.6}%
  \noindent\tabularx{\textwidth}{@{}>{\raggedright\arraybackslash}p{0.2\textwidth}X@{}}
  \arrayrulecolor{hkaRule}\hline
}{%
  \endtabularx
}
\newcommand{\kv}[2]{\leavevmode{\footnotesize\bfseries\color{hkaBlue}\MakeUppercase{#1}} & \leavevmode #2 \\ \hline}
\newcommand{\kvred}[2]{\leavevmode{\footnotesize\bfseries\color{hkaRed}\MakeUppercase{#1}} & \leavevmode #2 \\ \hline}

% --- Kleine Helfer --------------------------------------------------
\newcommand{\muted}[1]{{\color{hkaGray}#1}}
\newcommand{\highlight}[1]{{\color{hkaBlue}\bfseries#1}}
% Großer Kennzahl-Block: \metric{0.93}{ROC-AUC (val)}
\newcommand{\metric}[2]{%
  \begin{minipage}[t]{\linewidth}\centering
    {\fontsize{28}{32}\selectfont\bfseries\color{hkaBlue}#1\par}
    \vspace{1mm}{\footnotesize\color{hkaGray}#2\par}
  \end{minipage}}
% Bild-Platzhalter, bis ein echtes Bild da ist
\newcommand{\placeholder}[3][Figure]{%
  \begin{tikzpicture}
    \draw[hkaRule, fill=hkaStripe, rounded corners=2pt] (0,0) rectangle (#2,#3);
    \node[text=hkaLight, font=\small] at (#2/2,#3/2) {#1};
  \end{tikzpicture}}

% --- Code-Listings ---------------------------------------------------
\lstset{
  basicstyle=\ttfamily\scriptsize\color{hkaText},
  keywordstyle=\color{hkaBlue}\bfseries,
  commentstyle=\color{hkaLight}\itshape,
  stringstyle=\color{hkaRed},
  backgroundcolor=\color{hkaStripe},
  frame=none, framesep=2mm, xleftmargin=2mm, xrightmargin=2mm,
  showstringspaces=false, breaklines=true, columns=fullflexible
}

% =====================================================================
\begin{document}

% ---------------------------------------------------------------------
\begingroup
\setbeamertemplate{footline}{}
\begin{frame}[noframenumbering]
  \titlepage
\end{frame}
\endgroup

% ---------------------------------------------------------------------
\kicker{\week}
\begin{frame}{Agenda}
  \begin{enumerate}
    \item Recap: where we stood last week
    \item What we did this week
    \item Results
    \item Problems \& open questions
    \item Next steps
  \end{enumerate}
\end{frame}

% =====================================================================
\section{Recap}
% =====================================================================

\begin{frame}{Project at a glance}{Distinguish student-written essays from LLM-generated ones}
  \begin{kvtable}
    \kv{Dataset}{Kaggle \emph{LLM -- Detect AI Generated Text} (\texttt{train\_essays.csv})}
    \kv{Baseline}{TF-IDF + logistic regression; metric: ROC-AUC}
    \kvred{Experiment 1}{\muted{Placeholder: describe the experiment}}
    \kvred{Experiment 2}{\muted{Placeholder: describe the experiment}}
    \kv{Stretch goal}{\muted{Placeholder}}
    \kv{Pitfalls}{Very few AI-generated examples in the training data; risk of overfitting to prompts}
  \end{kvtable}
\end{frame}

% =====================================================================
\section{This Week}
% =====================================================================

\begin{frame}{What we did}{Short summary in one sentence}
  \begin{columns}[T,onlytextwidth]
    \column{0.48\textwidth}
      \textbf{Done}
      \begin{itemize}
        \item Data loading \& cleaning
        \item First baseline trained
        \item \alert{Important finding} highlighted
      \end{itemize}
    \column{0.48\textwidth}
      \textbf{In progress}
      \begin{itemize}
        \item Data augmentation with LLM-generated essays
        \item Cross-validation by prompt
          \begin{itemize}
            \item Sub-item
          \end{itemize}
      \end{itemize}
  \end{columns}
\end{frame}

\begin{frame}[fragile]{Pipeline}{Code snippets with \texttt{[fragile]}}
  \begin{columns}[T,onlytextwidth]
    \column{0.55\textwidth}
\begin{lstlisting}[language=Python]
from sklearn.feature_extraction.text import TfidfVectorizer
from sklearn.linear_model import LogisticRegression

# fit baseline
vec = TfidfVectorizer(ngram_range=(1, 2))
X = vec.fit_transform(df["text"])
clf = LogisticRegression().fit(X, df["generated"])
\end{lstlisting}
    \column{0.4\textwidth}
      \begin{exampleblock}{Note}
        Use \texttt{[fragile]} on every frame that contains code.
      \end{exampleblock}
  \end{columns}
\end{frame}

% =====================================================================
\section{Results}
% =====================================================================

\begin{frame}{Key numbers}
  \vspace{6mm}
  \begin{columns}[T,onlytextwidth]
    \column{0.32\textwidth}\metric{0.00}{ROC-AUC (validation)}
    \column{0.32\textwidth}\metric{0.00}{Accuracy}
    \column{0.32\textwidth}\metric{0}{Training essays}
  \end{columns}
  \vspace{10mm}
  \begin{center}
    \muted{One-sentence takeaway: what do these numbers mean?}
  \end{center}
\end{frame}

\begin{frame}{Figure with explanation}
  \begin{columns}[c,onlytextwidth]
    \column{0.58\textwidth}
      % \includegraphics[width=\linewidth]{figures/roc.pdf}
      \placeholder[ROC curve]{7.6}{4.6}
    \column{0.38\textwidth}
      \begin{itemize}
        \item Observation 1
        \item Observation 2
        \item \alert{What surprised us}
      \end{itemize}
  \end{columns}
\end{frame}

\begin{frame}{Model comparison}
  \centering
  \renewcommand{\arraystretch}{1.3}
  \begin{tabular}{@{}lccc@{}}
    \toprule
    \textbf{Model} & \textbf{ROC-AUC} & \textbf{Accuracy} & \textbf{Train time} \\
    \midrule
    TF-IDF + LogReg & 0.00 & 0.00 & 0\,s \\
    \rowcolor{hkaStripe}
    DeBERTa-v3 (fine-tuned) & \highlight{0.00} & 0.00 & 0\,min \\
    Ensemble & 0.00 & 0.00 & 0\,min \\
    \bottomrule
  \end{tabular}
\end{frame}

% =====================================================================
\section{Open Questions}
% =====================================================================

\begin{frame}{Problems \& questions for the supervisor}
  \begin{alertblock}{Problem}
    Describe a blocker in one or two sentences.
  \end{alertblock}
  \vspace{3mm}
  \begin{block}{Question}
    What we would like feedback on.
  \end{block}
\end{frame}

\kicker{Outlook}
\begin{frame}{Next steps}
  \begin{kvtable}
    \kv{Next week}{\muted{Concrete goal for the coming week}}
    \kv{Who}{\muted{Who does what}}
    \kvred{Risk}{\muted{What could go wrong}}
  \end{kvtable}
\end{frame}

% ---------------------------------------------------------------------
\begingroup
\setbeamertemplate{footline}{}
\begin{frame}[noframenumbering]
  \centering\vfill
  {\fontsize{36}{42}\selectfont\bfseries\color{hkaBlueTitle}Thank you!\par}
  \vspace{4mm}
  {\large Questions?\par}
  \vspace{8mm}
  \hkalogo\par
  \vfill
\end{frame}
\endgroup

\end{document}
